% !TEX root = ../main.tex
%-----------------------------------------------------------------------
%  What it costs, the gradient-check saving, and the request.
%  \input by ../main.tex -- not compiled on its own.
%-----------------------------------------------------------------------
\hd{What it costs.}
Both halves of a member have been measured on completed runs at this
configuration and rank count, so these are not extrapolations. Two adjoint
measurements at very different window lengths agree to $3\,\%$, so cost is
linear in window length.

\smallskip
\begin{tabularx}{\textwidth}{@{}Y{1.6}Y{0.6}Y{0.8}@{}}
\toprule
\textbf{Configuration} & \textbf{Per member} & \textbf{Seven members}\\
\midrule
As the reference run was configured & 24.7 h & \textbf{173 node-h}\\
With the built-in gradient check disabled & $\sim$10.5 h & \textbf{$\sim$75 node-h}\\
Gradient check disabled, four runs packed per node & --- & \textbf{$\sim$19 node-h}\\
\bottomrule
\end{tabularx}
\smallskip

\hd{Most of that cost is avoidable, and should be avoided.}
Every adjoint job in this configuration also runs the model's built-in
finite-difference gradient check, and pays for the perturbed forward
integrations it needs. On an instrumented 30-day run the model executed
$18{,}622$ forward steps in order to integrate $1{,}440$ --- a factor of
$12.9$, and the reason the adjoint costs about $28$ forward steps per timestep
where checkpointing alone would predict three to five. Strip the check out and
the remainder falls squarely inside that expected range.

Turning it off for ensemble members is the largest single saving available and
the easiest to take. Nothing is lost by doing so, for the reason the next
section gives.

\hd{The request.}

\smallskip
\begin{tabularx}{\textwidth}{@{}Y{0.42}Y{1.58}@{}}
\toprule
\textbf{Compute} & \textbf{$\sim$75 node-hours} expected; \textbf{200 requested}
  as a ceiling\\
\textbf{Node type} & CPU only, MPI, no GPU. 27 ranks per run\\
\textbf{Storage} & $\sim$23 GB retained. Not a constraint\\
\textbf{Duration} & 2--3 weeks including setup and analysis\\
\bottomrule
\end{tabularx}
\smallskip

The ceiling is requested rather than the estimate for three reasons: the
$2.6\times$ saving above is inferred from call counts rather than measured on a
completed run; strongly perturbed members may need a reduced timestep or extra
solver iterations; and some of seven never-before-integrated settings will need
repeating. A benchmark costing well under one node-hour settles all three, and
it is the first step of the sequence below.
